\documentclass[aspectratio=169,10pt]{beamer}
\usepackage[T1]{fontenc}
\usepackage{lmodern,graphicx,booktabs,tabularx,xcolor}
\definecolor{lhcbgreen}{HTML}{16836B}
\definecolor{lhcbink}{HTML}{20303C}
\setbeamercolor{structure}{fg=lhcbgreen}
\setbeamercolor{normal text}{fg=lhcbink,bg=white}
\setbeamercolor{frametitle}{fg=lhcbink,bg=white}
\setbeamertemplate{navigation symbols}{}
\setbeamertemplate{footline}[frame number]
\setbeamersize{text margin left=8mm,text margin right=8mm}
\setlength{\tabcolsep}{7pt}
\renewcommand{\arraystretch}{1.2}
\newcommand{\asset}[1]{assets/transport_ft_with_hole/#1}
\newcommand{\HoleWidth}{3288.4}
\newcommand{\HoleHeight}{672.6}
\newcommand{\HoleArea}{1.3806}
\newcommand{\MPArea}{1.3003}
\newcommand{\GapArea}{0.0803}
\newcommand{\FieldIntegral}{-4.1856}

\newcommand{\InventoryRows}{
TV & 50 & 50 & -287.5--557.6 \\
RICH1 & 2 & 2 & 795.0--2237.2 \\
UP & 16 & 4 & 2314.0--2656.0 \\
MP & 5 & 5 & 7665.5--9132.5 \\
FT & 12 & 12 & 7902.7--9384.7 \\
RICH2 & 2 & 2 & 9487.5--11839.2 \\
PicoCal & 2 & 2 & 12565.9--12601.0 \\
MUON & 4 & 4 & 15270.0--18870.0 \\
}

\newcommand{\AreaRows}{
8 & 18.3690 & 17.0560 & 7.15\% \\
10 & 22.9781 & 21.6652 & 5.71\% \\
12 & 27.5873 & 26.2743 & 4.76\% \\
}

\newcommand{\ToyRows}{
Neutral control & 50 & 5 & 12 & 0 & 4 & 4 \\
Positive & 50 & 3 & 12 & 7 & 4 & 4 \\
Negative & 50 & 5 & 12 & 0 & 4 & 4 \\
Positive & 10 & 0 & 12 & 12 & 4 & 4 \\
Negative & 10 & 5 & 12 & 0 & 4 & 4 \\
}

\title{LHCb Detector Masks and Field Transport}
\subtitle{Configurable XY surfaces and the FT\_withHole variants}
\date{18 October 2026}

\begin{document}
\begin{frame}
\titlepage
\begin{center}\small Six original setups + six FT\_withHole variants\\
TV, RICH1, UP, MP, FT, RICH2, PicoCal, and Muon\end{center}
\end{frame}

\begin{frame}{Detector planes and logical layers}
\begin{columns}[T,onlytextwidth]
\column{0.54\textwidth}
\small
\setlength{\tabcolsep}{3pt}
\begin{tabular}{@{}lrrr@{}}
\toprule
System & Surfaces & Layers & Global $z$ range [mm]\\
\midrule
\InventoryRows
\bottomrule
\end{tabular}
\column{0.41\textwidth}
\small Each physical surface has a global $z$, an XY polygon, and a stable
logical layer label. A logical layer fires if any of its physical surfaces
contains the transported crossing.\\[3mm]
UP keeps 16 physical planes and counts four logical layers. FT keeps
X/U/V/X in each of T1, T2, and T3: 12 logical layers.\\[3mm]
FT\_withHole keeps every plane position and label. Only the FT XY masks change.
\end{columns}
\end{frame}

\begin{frame}{Global Z placement of the detector surfaces}
\centering\includegraphics[width=\textwidth,height=0.68\textheight,keepaspectratio]{\asset{z_layout.pdf}}
\par\vspace{3mm}\small Every tick is a physical surface Z position. UP contains
four groups of four surfaces; each FT station contains X/U/V/X.\\[2mm]
The FT\_withHole variant changes the XY masks while retaining all placements.
\end{frame}

\begin{frame}{TV, RICH1, UP, and MP XY acceptance}
\centering\includegraphics[width=\textwidth,height=0.78\textheight,keepaspectratio]{\asset{xy_inventory_1.pdf}}
\par\scriptsize XY unions of stored surfaces; axes are scaled independently.
Open sensor gaps are shown as gaps rather than assigned a rectangular hole size.
\end{frame}

\begin{frame}{FT, RICH2, PicoCal, and Muon XY acceptance}
\centering\includegraphics[width=\textwidth,height=0.78\textheight,keepaspectratio]{\asset{xy_inventory_2.pdf}}
\par\scriptsize XY unions of stored surfaces. FT: nominal outline in blue,
retained FT\_withHole in green, stepped MP envelope in red.
\end{frame}

\begin{frame}{The three UP border variants}
\centering\includegraphics[width=\textwidth,height=0.74\textheight,keepaspectratio]{\asset{up_variants.pdf}}
\small
\begin{tabularx}{\textwidth}{@{}lX@{}}
\texttt{UP\_GDML} & Original extracted acceptance.\\
\texttt{UP\_Y8} & Top/bottom cutoff.\\
\texttt{UP\_Y8X6} & Top/bottom and left/right cutoffs.\\
\end{tabularx}
\end{frame}

\begin{frame}{FT\_withHole joins the MP sides near X=0}
\centering
$H_{\rm MP}=\operatorname{join}_{x=0}\!\left(\bigcup_{i=1}^{5} A_{{\rm MP},i}\right),
\qquad A_{{\rm FT},j}^{\rm withHole}=A_{{\rm FT},j}\setminus H_{\rm MP}$
\par\vspace{2mm}
\includegraphics[width=\textwidth,height=0.61\textheight,keepaspectratio]{\asset{mp_hole_construction.pdf}}
\par\small
Hole envelope: \HoleWidth{} $\times$ \HoleHeight{} mm; area \HoleArea{} m$^2$.\\
MP active union: \MPArea{} m$^2$; the envelope additionally excludes
\GapArea{} m$^2$ across the inner gap; the outer square steps are preserved.\\[1mm]
One global XY envelope is used at every FT plane; there is no projection in $z$.
\end{frame}

\begin{frame}{FT 8 modules: original and FT\_withHole}
\centering\includegraphics[width=\textwidth,height=0.86\textheight,keepaspectratio]{\asset{ft_compare_8.pdf}}
\end{frame}

\begin{frame}{FT 10 modules: original and FT\_withHole}
\centering\includegraphics[width=\textwidth,height=0.86\textheight,keepaspectratio]{\asset{ft_compare_10.pdf}}
\end{frame}

\begin{frame}{FT 12 modules: original and FT\_withHole}
\centering\includegraphics[width=\textwidth,height=0.86\textheight,keepaspectratio]{\asset{ft_compare_12.pdf}}
\end{frame}

\begin{frame}{Setup inventory and FT active area}
\begin{columns}[T,onlytextwidth]
\column{0.5\textwidth}
\small
\begin{tabular}{@{}lccc@{}}
\toprule
UP choice & FT modules & Original & With hole\\
\midrule
GDML & 8/8/10 & Yes & Yes\\
Y8 & 8/8/10 & Yes & Yes\\
Y8X6 & 8/8/10 & Yes & Yes\\
GDML & 10/10/12 & Yes & Yes\\
Y8 & 10/10/12 & Yes & Yes\\
Y8X6 & 10/10/12 & Yes & Yes\\
\bottomrule
\end{tabular}\\[3mm]
Each original catalog ID has an additional ID ending in
\texttt{\_FT\_withHole}. Other detector masks are identical to its parent setup.
\column{0.46\textwidth}
\small
\begin{tabular}{@{}rrrr@{}}
\toprule
Modules & Original & Kept & Removed\\
 & [m$^2$] & [m$^2$] & \\
\midrule
\AreaRows
\bottomrule
\end{tabular}\\[3mm]
Areas are per FT plane. X/U/V rotations preserve these areas for the current
MP envelope. The original area already excludes the 260 mm beam opening.
\end{columns}
\end{frame}

\begin{frame}{Field map and magnet polarity}
\centering\includegraphics[width=\textwidth,height=0.60\textheight,keepaspectratio]{\asset{field_polarity.pdf}}
\small
The ROOT service reads Bx, By, Bz on a 100 mm grid. Coordinates are mm,
field is tesla, momentum is MeV/$c$.\\[2mm]
\texttt{polarity=-1} reverses all three components, including fringe fields.\\
On-axis $\int_0^{10000\,{\rm mm}} B_y\,dz=\FieldIntegral$ T\,m for the supplied map.
\end{frame}

\begin{frame}{Toy trajectories through the full detector}
\centering\includegraphics[width=\textwidth,height=0.65\textheight,keepaspectratio]{\asset{toy_trajectories.pdf}}
\small
Initial slopes $t_x=0.06$, $t_y=0.03$; start at $z=-300$ mm with the zero-field
line through the origin. The same field trajectory is evaluated against both mask variants.\\[2mm]
RK4 transport in $z$; cubic Hermite XY interpolation. No material effects.
A nearly transverse or curling path raises \texttt{TrajectoryError}.
\end{frame}

\begin{frame}{Fired layers before and after the enlarged hole}
\small Reference setup: UP Y8X6, FT 8/8/10; supplied field polarity $+1$.
\par\vspace{3mm}\centering
\begin{tabular}{@{}lrrrrrr@{}}
\toprule
Track & $p$ [GeV/$c$] & nMP & nFT original & nFT with hole & nUP & nMuon\\
\midrule
\ToyRows
\bottomrule
\end{tabular}
\par\vspace{4mm}\raggedright
FT\_withHole can only reduce FT acceptance. MP, UP, TV, RICH, PicoCal, and
Muon use the same masks and return the same flags as the parent setup.\\[3mm]
An XY point in an MP sensor gap is also rejected by FT when it lies inside
the solid envelope. Removing FT overlap does not guarantee an MP hit at a different $z$.
\end{frame}

\begin{frame}{Batch and parallel propagation}
\begin{tabular}{@{}lrr@{}}
\toprule
Method at 10 mm steps & Per track & 4 million tracks\\
\midrule
Individual Z spline & 314 ms & 14.6 days\\
Batch of 1024, one process & 1.25 ms & 1.4 hours\\
Batch of 1024, four processes & 0.45 ms & 30 minutes\\
\bottomrule
\end{tabular}
\par\vspace{4mm}\small
Measured on 18 October 2026. Batch timing: 8192 toy tracks at 10--100 GeV/$c$,
including labels; four-process timing includes worker startup. Four-million
times are extrapolations. Serial/parallel hit-count checksums agreed.\\[3mm]
The individual Z path uses about 2005 nodes; vectorized batches use roughly
1500 main steps with longer drift steps outside the field. Reuse crossings to
compare geometries sharing the same physical Z planes.\\[3mm]
50 mm steps reduced cost further: maximum XY difference 0.0272 mm, with no changed
layer flags on a 4096-track check. The default remains 10 mm.
\end{frame}

\begin{frame}[fragile]{Batch mask columns and geometry comparisons}
\small
\begin{verbatim}
columns = dict(x=x, y=y, z=z, tx=tx, ty=ty, q_over_p=qop)
batch = transport.propagate_batch(columns, label_suffix="base")
track_table.update(batch.as_columns())

# Different XY mask, same state/field/Z planes: reuse transport.
other = DetectorTransport(other_configuration, transport.field)
track_table.update(other.labels_from_crossings(
    batch.crossings, label_suffix="withHole",
))
\end{verbatim}
\vspace{3mm}
\small A worker loads the geometry and field once, then processes chunks of
1024--4096 tracks. The process-pool example is
\texttt{examples/benchmark\_transport.py}.\\[2mm]
Batch transport raises on failed/incomplete tracks, reporting their row indices.
\end{frame}

\begin{frame}[fragile]{Loading a setup and appending track flags}
\small
\begin{verbatim}
from xyshapes_detectors import DetectorTransport, TrackState

transport = DetectorTransport.from_files(
    "lhcb_lower_up_y8x6_ft_8_8_10_FT_withHole",
    "MagField/magnetic_field_v5r0.root", polarity=+1,
)
state = TrackState.from_charge_momentum(
    charge=1, momentum_mev=50_000,
    x=0, y=0, tx=0.06, ty=0.03, z=-300,
)
result = transport.propagate(state, label_suffix="hole")
track_row.update(result.as_record())
\end{verbatim}
\small
\texttt{result.layer\_hits}: explicit logical-layer booleans.\\
\texttt{result.trajectory}: path and XY intercepts at requested $z$.\\
Returned fields include \texttt{nTV\_hole}, \texttt{nUP\_hole}, \texttt{nMP\_hole},
\texttt{nFT\_hole}, \texttt{nMuon\_hole}, RICH flags, and \texttt{hasPicoCal\_hole}.
\end{frame}

\begin{frame}{Geometry derivation and reusable files}
\small
\begin{enumerate}
\item Detector surfaces and placements originate from the saved ViewLHCb/GDML
acceptances and simplified FT module-count constructions.
\item \texttt{build\_ft\_with\_hole\_variants.py} joins the two MP sides near X=0
within each horizontal band, subtracts this stepped hole from each FT
mask, and saves named shape pickles.
\item Every new setup stores \texttt{layout.csv}, \texttt{config.json}, and
\texttt{ft\_area\_summary.json}. The config records the envelope vertices and operation.
\item Catalog entries reuse shared source masks and nine derived FT shapes.
Both catalog IDs and setup directory paths can be loaded.
\end{enumerate}
\vspace{3mm}
\textbf{Sources:} \texttt{detector\_setups/catalog.json}, saved XY pickles,
\texttt{variants.py}, \texttt{transport.py}, and the supplied field ROOT file.\\[2mm]
\textbf{Scope:} geometric acceptance along a vacuum trajectory. This is not a
material, scattering, light-transport, or detector-efficiency model.
\end{frame}
\end{document}
